% Original PDF 29 / printed 19. All of Box 2.1, not included in accepted sample.
% Source box is represented using the existing user's template color, not a new style.
% block: homological-algebra-summary-heading page:29 type:ordinary-paragraph
\begin{tcolorbox}[enhanced,breakable,colback=KuangNei,frame hidden,boxrule=0pt,sharpish corners,before skip=10pt,after skip=10pt]
\refstepcounter{booksummary}\label{box:homological-algebra-summary}
\noindent\textbf{Box~\thebooksummary:} {Summary of Chapter~\ref{chap:homological-algebra} on homological algebra.}\par

% block: summary-chain-complex-intro page:29 type:ordinary-paragraph
{A \emph{chain complex} is a sequence of vector spaces and operators}

% block: summary-chain-sequence page:29 type:equation
\[
\cdots\to V_{k+1}\xrightarrow{\partial_{k+1}}V_k\xrightarrow{\partial_k}V_{k-1}\to\cdots
\]

% block: summary-chain-spaces page:29 type:ordinary-paragraph
{with successive operators composing to zero. It gives rise to three vector spaces at each level $k$:}
\begin{itemize}
\item {$\mathcal Z_k$, the null space of the boundary operator $\partial_k$, consisting of the \emph{cycles};}
\item {$\mathcal B_k$, the range of $\partial_{k+1}$, the \emph{boundaries};}
\item {$\mathcal H_k=\mathcal Z_k/\mathcal B_k$ the \emph{homology space}.}
\end{itemize}

% block: summary-simplicial-example page:29 type:ordinary-paragraph
{An important example is the \emph{simplicial chain complex} associated to a triangulation or more general simplicial complex. The dimensions of its homology spaces are the Betti numbers of the underlying space.}

% block: summary-cochain-intro page:29 type:ordinary-paragraph
{A \emph{cochain complex} has the indexing reversed,}

% block: summary-cochain-sequence page:29 type:equation
\[
\cdots\to V^{k-1}\xrightarrow{d^{k-1}}V^k\xrightarrow{d^k}V^{k+1}\to\cdots,
\]

% block: summary-cochain-objects page:29 type:ordinary-paragraph
{giving rise to \emph{cocyles}, \emph{coboundaries}, and \emph{cohomology}.}

% block: summary-duality page:29 type:ordinary-paragraph
{The dual spaces and dual maps of a chain complex form a cochain complex.}

% block: summary-de-rham-intro page:29 type:ordinary-paragraph
{Another cochain complex is the \emph{de Rham complex} which, on a three-dimensional domain, can be written}

% block: summary-de-rham-sequence page:29 type:equation
\[
0\to C^\infty(\Omega)\xrightarrow{\operatorname{grad}}C^\infty(\Omega;\mathbb R^3)\xrightarrow{\operatorname{curl}}C^\infty(\Omega;\mathbb R^3)\xrightarrow{\operatorname{div}}C^\infty(\Omega)\to0.
\]

% block: summary-de-rham-theorem page:29 type:ordinary-paragraph
{The de Rham map is a \emph{cochain map} from the de Rham complex to the dual of the simplicial chain complex and, as such, induces a map on cohomology. De Rham's theorem states that the induced map is an isomorphism.}
\end{tcolorbox}
